\begin{dang}{Bảo toàn năng lượng và bảo toàn động lượng}	
\Anlythuyet{
\begin{itemize}
\item $\begin{cases}
\ct{Động năng: } & \ K=\dfrac{mv^2}{2}\\ 
\ct{Động lượng: } &\vec p=m\vec v
\end{cases}
\Rightarrow \boder{p^2=2mK}$	
\item Áp dụng định luật bảo toàn năng lượng toàn phần:
\begin{align*}
\varepsilon_\text{trước}+m_\text{trước}.c^2+K_{\text{trước}}=\varepsilon_\text{sau}+m_\text{sau}.c^2+K_{sau}\\			\Rightarrow\ \boder{\varepsilon_\text{trước}+W+K_{\text{trước}}=\varepsilon_\text{sau}+K_{sau}}\ (1)
\end{align*}
Với $\varepsilon=hf$ là năng lượng của một photon bức xạ.
\item Áp dụng định luật bảo toàn động lượng, vận dụng quy tắc hình bình hành và các định lí trong tam giác suy ra công thức dưới dạng độ lớn.
\end{itemize}\begin{itemize}
\item Xét phản ứng dùng hạt nhân A làm "đạn" bắn vào hạt nhân B là "bia" đứng yên, sinh ra hai hạt nhân C và D và phản ứng không kèm theo bức xạ $\gamma$.
\begin{itemize}
\item Bảo toàn năng lượng: \boder{W =K_\text{sau}-K_\text{trước}}\quad $(2)$
\item Bảo toàn động lượng: \boder{\vec{p}_A=\vec{p}_C+\vec{p}_D}\quad $(3)$
\item Động năng tối thiểu của hạt "đạn" A để phản ứng xảy ra là: $\boder{K_{A\min}=|W|}$
\end{itemize}
\item Nếu dùng bức xạ $\gamma$ để thực hiện phản ứng hạt nhân thì mỗi hạt nhân chỉ hấp thụ một photon\\
Năng lượng tối thiểu của bức xạ $\gamma$ để phản ứng xảy ra: \boder{\varepsilon_{\min}=hf_{\min}=\dfrac{hc}{\lambda_{\max}}=|W|}
\end{itemize}
}
\end{dang}
